\section{Implementación Práctica del Entorno HTTPS con Certificados Digitales}

La presente sección documenta de manera formal y detallada la implementación práctica desarrollada en el marco de la asignatura de Criptografía de la Facultad de Ingeniería de la Universidad Central del Ecuador. El objetivo del ejercicio es construir un entorno seguro basado en un servidor HTTPS utilizando certificados digitales X.509 generados con OpenSSL, integrados dentro de una arquitectura modular implementada en Node.js, Express.js y una interfaz gráfica web desarrollada con tecnologías HTML, CSS y JavaScript.

Se expone la estructura completa del proyecto, la generación profesional de certificados digitales, su proceso de verificación, así como la integración con el backend y el frontend. Finalmente, se reservan subsecciones específicas para registrar el código fuente íntegro asociado a la solución, con el propósito de mantener trazabilidad técnica y evidencia documental del desarrollo.

\subsection{Estructura General del Proyecto}

El proyecto se organiza bajo una arquitectura modular, permitiendo separar la capa de servidor, la capa de cliente y los componentes de seguridad. La estructura completa es la siguiente:

\begin{verbatim}
secure-web/
├── certs/                  # Certificados y claves criptográficas
│   ├── server.key          # Clave privada del servidor
│   └── server.crt          # Certificado digital X.509 autofirmado
├── public/                 # Interfaz gráfica y archivos estáticos
│   ├── css/
│   ├── js/
│   │   └── main.js
│   └── index.html
├── src/                    # Backend implementado en Node.js
│   ├── app.js
│   └── server.js
├── package.json            # Dependencias y metadatos del proyecto
└── README.md               # Documentación técnica
\end{verbatim}

Esta estructura facilita el análisis del flujo de información, el aislamiento de responsabilidades y la comprensión del proceso de comunicación cifrada mediante TLS.

\subsection{Tecnologías Empleadas}

En la Tabla \ref{tabla:tecnologias} se presentan las principales tecnologías utilizadas en el proyecto:

\begin{table}[H]
\centering
\begin{tabularx}{\textwidth}{lX}
\textbf{Tecnología} & \textbf{Descripción} \\
\hline
Node.js & Entorno de ejecución para JavaScript del lado del servidor. \\
Express.js & Framework minimalista para la creación de APIs y servidores HTTP/HTTPS. \\
OpenSSL & Herramienta estándar para la generación y gestión de certificados X.509. \\
HTML5 / CSS3 & Tecnologías base para el desarrollo de la interfaz gráfica. \\
JavaScript ES6+ & Lógica dinámica en cliente y servidor. \\
Módulo Crypto (Node.js) & Extracción de huellas digitales y análisis criptográfico. \\
\end{tabularx}
\caption{Tecnologías empleadas en el desarrollo del proyecto.}
\label{tabla:tecnologias}
\end{table}

\subsection{Descripción de Componentes Principales}

\subsubsection*{Backend (Directorio \texttt{src/})}

\begin{itemize}
    \item \textbf{server.js}: Punto de entrada del servidor HTTPS.  
    Carga los certificados digitales, configura las opciones TLS y levanta el servidor en el puerto 3000.
    \item \textbf{app.js}: Define la configuración principal de Express, el enrutamiento, el middleware y el endpoint encargado de analizar y exponer la información contenida en el certificado X.509.
\end{itemize}

\subsubsection*{Frontend (Directorio \texttt{public/})}

\begin{itemize}
    \item \textbf{index.html}: Interfaz visual donde se muestra la información del certificado.
    \item \textbf{main.js}: Lógica del cliente encargada de consumir la API REST, actualizar el DOM, mostrar los valores criptográficos y gestionar la experiencia interactiva.
\end{itemize}

\subsubsection*{Seguridad (Directorio \texttt{certs/})}

\begin{itemize}
    \item \textbf{server.key}: Clave privada RSA utilizada en el handshake TLS.
    \item \textbf{server.crt}: Certificado digital que se entrega a los clientes para autenticar la identidad del servidor.
\end{itemize}

\subsection{Generación de Certificados SSL/TLS utilizando Git Bash}

La creación de certificados digitales constituye un procedimiento fundamental dentro del estudio de la seguridad informática y, específicamente, en el contexto de la cátedra de Criptografía de la Universidad Central del Ecuador (UCE). Con el fin de implementar un entorno seguro para servicios web basados en HTTPS, se empleó la herramienta \textit{OpenSSL} ejecutada desde el entorno \textit{Git Bash} en sistemas Windows, debido a su compatibilidad, facilidad de uso y amplio soporte para operaciones criptográficas.

Este proceso permite generar tanto la clave privada del servidor como un certificado digital autofirmado bajo el estándar X.509, el cual será posteriormente utilizado por un servidor HTTPS implementado en Node.js. El objetivo académico es comprender la estructura del certificado, sus campos internos, su finalidad dentro de un sistema criptográfico y la mecánica del proceso de firma digital.

\subsubsection{Creación del directorio de certificados}

El primer paso consiste en la creación de un directorio destinado exclusivamente al almacenamiento seguro de los archivos criptográficos:

\begin{verbatim}
mkdir certs
cd certs
\end{verbatim}

Este directorio contendrá dos archivos esenciales para la configuración del servidor HTTPS:

\begin{itemize}
    \item \textbf{server.key}: clave privada RSA generada durante el proceso.
    \item \textbf{server.crt}: certificado digital autofirmado.
\end{itemize}

\subsubsection{Generación del certificado y la clave privada}

Desde \textit{Git Bash}, se ejecuta el siguiente comando, el cual genera automáticamente una clave privada RSA de 2048 bits y un certificado digital autofirmado con una validez de un año:

\begin{verbatim}
openssl req -x509 -nodes -newkey rsa:2048 \
-keyout server.key -out server.crt -days 365
\end{verbatim}

Durante el proceso, OpenSSL solicitará información relevante para construir la estructura X.509. Para fines académicos y manteniendo coherencia institucional, se utilizaron los siguientes valores recomendados:

\begin{verbatim}
Country Name (2 letter code) [AU]: EC
State or Province Name (full name): Pichincha
Locality Name (eg, city): Quito
Organization Name (eg, company): Universidad Central del Ecuador
Organizational Unit Name (eg, section): Facultad de Ingeniería,
                                        Escuela de Computación
Common Name (e.g. server FQDN or YOUR name): localhost
Email Address: criptografia@uce.edu.ec
\end{verbatim}

Es importante destacar que el campo \textbf{Common Name (CN)} debe establecerse como \textit{localhost}, dado que el proyecto se ejecuta en un entorno de desarrollo local.

\subsubsection{Verificación del certificado generado}

Una vez creado el certificado, es posible inspeccionar su contenido utilizando el comando:

\begin{verbatim}
openssl x509 -in server.crt -text -noout
\end{verbatim}

Esta instrucción permite revisar:

\begin{itemize}
    \item Los campos del sujeto (Subject) y emisor (Issuer).
    \item La clave pública RSA asociada.
    \item El algoritmo de firma digital empleado.
    \item El período de validez.
    \item Las huellas digitales (\textit{fingerprints}) SHA--1 y SHA--256.
    \item Las extensiones X.509 del certificado.
\end{itemize}

\subsubsection{Regeneración del certificado}

En caso de requerir un nuevo certificado, se puede repetir el proceso eliminando los archivos anteriores:

\begin{verbatim}
cd certs
rm server.key server.crt
openssl req -x509 -nodes -newkey rsa:2048 \
-keyout server.key -out server.crt -days 365
\end{verbatim}

\subsubsection{Integración con el servidor HTTPS}

El certificado y la clave privada generados desde Git Bash se integran directamente en el servidor implementado en Node.js mediante:

\begin{verbatim}
const httpsOptions = {
  key: fs.readFileSync('./certs/server.key'),
  cert: fs.readFileSync('./certs/server.crt'),
};
\end{verbatim}

De esta manera, el servidor establece conexiones seguras mediante el protocolo HTTPS, garantizando confidencialidad, integridad y autenticación básica del servidor en un entorno académico controlado.

\subsection{Integración con el Servidor HTTPS}

El servidor utiliza los certificados generados mediante la siguiente configuración:

\begin{verbatim}
const httpsOptions = {
  key: fs.readFileSync('./certs/server.key'),
  cert: fs.readFileSync('./certs/server.crt')
};
\end{verbatim}

Al actualizar o regenerar los certificados, la interfaz web reflejará automáticamente los nuevos valores al invocar el endpoint:

\begin{verbatim}
fetch('/api/cert-info')
\end{verbatim}

\section{Evidencia del Código Fuente}

Con el propósito de mantener una documentación completa y verificable, las siguientes subsecciones están destinadas a incluir el \textbf{código íntegro} del proyecto.

\subsection{Archivo \texttt{server.js}}

\begin{lstlisting}[language=JavaScript, caption={Implementación del servidor HTTPS en Node.js}, label={lst:serverjs}]
import https from "https";
import fs from "fs";
import path from "path";
import app from "./app.js";
import { fileURLToPath } from "url";

// Resolver la ruta del archivo actual
const __filename = fileURLToPath(import.meta.url);
const __dirname = path.dirname(__filename);

// Cargar certificados SSL/TLS
const key = fs.readFileSync(path.join(__dirname, "../certs/server.key"));
const cert = fs.readFileSync(path.join(__dirname, "../certs/server.crt"));

// Crear el servidor HTTPS
const httpsServer = https.createServer({ key, cert }, app);

// Puerto de escucha
const PORT = 3000;

// Iniciar el servidor
httpsServer.listen(PORT, () => {
  console.log(`Servidor HTTPS ejecutándose en https://localhost:${PORT}`);
});
\end{lstlisting}

\subsubsection*{Descripción detallada del archivo \texttt{server.js}}
\label{subsubsec:serverjs_descripcion}

El código mostrado en el \textit{Listing}~\ref{lst:serverjs} implementa el arranque de un servidor HTTPS en Node.js, integrando certificados X.509 generados con OpenSSL y una aplicación Express exportada desde \texttt{app.js}. A continuación se desglosa el propósito y el comportamiento de cada bloque funcional del fichero, con recomendaciones operacionales y de seguridad.

\paragraph{1. Importación de módulos}
\begin{description}
  \item[\texttt{https}] Módulo nativo de Node.js que permite crear servidores TLS/SSL. Se utiliza para levantar un servidor que escucha conexiones cifradas y realiza el \textit{handshake} TLS con los clientes.
  \item[\texttt{fs}] Módulo de sistema de archivos, empleado para leer los artefactos criptográficos (clave privada y certificado) desde disco.
  \item[\texttt{path}] Utilizado para construir rutas de archivos portables entre distintos sistemas operativos (evita concatenaciones de cadenas inseguras).
  \item[\texttt{app}] Exportación de la instancia Express configurada (rutas, middleware y lógica de la aplicación) — definida en \texttt{src/app.js}.
  \item[\texttt{fileURLToPath}] Función de la librería \texttt{url} para transformar la URL del módulo (ESM) en una ruta de sistema de archivos; se utiliza cuando el proyecto usa módulos ECMAScript nativos (``type": "module'').
\end{description}

\paragraph{2. Resolución de la ruta del fichero actual}
Las constantes \texttt{\_\_filename} y \texttt{\_\_dirname} se obtienen mediante las funciones \texttt{fileURLToPath(import.meta.url)} y \texttt{path.dirname(...)}. Dado que el proyecto utiliza módulos ES (ECMAScript Modules — ESM) en lugar del sistema clásico CommonJS, estas funciones son necesarias para resolver correctamente las rutas absolutas del proyecto. Esto permite acceder sin ambigüedades a archivos relativos —por ejemplo \texttt{../certs/server.key}— independientemente del directorio desde el que se ejecute el servidor o del sistema operativo utilizado. De esta forma se mitigan errores comunes en tiempo de ejecución asociados a rutas mal resueltas.

\paragraph{3. Carga de los certificados SSL/TLS}
El servidor HTTPS carga los certificados mediante \texttt{fs.readFileSync(path.join(\_\_dirname, "../certs/server.key"))} y la lectura equivalente del certificado público. Durante esta fase se incorporan en memoria la \textbf{clave privada} y el \textbf{certificado digital}, elementos indispensables para establecer comunicaciones cifradas bajo TLS. Es importante considerar los siguientes aspectos:

\begin{itemize}
  \item \textbf{Protección de la clave privada:} El archivo \texttt{server.key} debe poseer permisos estrictos (por ejemplo, 600) y nunca debe almacenarse en repositorios públicos. En entornos de producción se recomienda emplear módulos HSM o servicios de gestión de claves (KMS) para evitar que la clave privada resida en disco.
  
  \item \textbf{Lectura síncrona:} El uso de \texttt{readFileSync} es apropiado durante la fase de arranque del servidor, cuando todavía no se aceptan conexiones. Para sistemas que requieren \textit{hot reload} de certificados o rotación periódica de claves, conviene implementar mecanismos de recarga segura o utilizar APIs que soporten rotación de credenciales sin interrumpir el servicio.
  
  \item \textbf{Validación previa:} Es recomendable verificar el contenido del certificado y de la clave mediante operaciones de parseo, validación de formato o detección de corrupción. Asimismo, la implementación de manejo de excepciones permite capturar de manera controlada errores derivados de rutas inválidas o archivos mal formados.
\end{itemize}


\paragraph{4. Creación del servidor HTTPS}
La llamada \texttt{https.createServer(\{ key, cert \}, app)} construye un servidor TLS que, una vez aceptada la conexión y completado el handshake, delega las peticiones HTTP en la aplicación Express (\texttt{app}). Aspectos relevantes:
\begin{itemize}
  \item \textbf{Opciones adicionales:} En un despliegue real conviene especificar opciones adicionales como \texttt{ca} (certificados intermedios), \texttt{honorCipherOrder}, \texttt{ciphers} y las curvas elípticas preferidas para controlar la política criptográfica del servidor.
  \item \textbf{Cipher suites y políticas:} Configurar cipher suites priorizando ECDHE+AEAD (por ejemplo AES-GCM o ChaCha20-Poly1305) y deshabilitar suites obsoletas según recomendaciones RFC y guías operacionales \cite{rfc8446, rfc7525}.
  \item \textbf{Integración con Express:} Express actúa como manejador de rutas y middleware; la separación entre la capa TLS y la lógica de aplicación facilita aplicar políticas de seguridad en ambos niveles.
\end{itemize}

\paragraph{5. Configuración de puerto y puesta en escucha}
Se define la constante \texttt{PORT} (3000 en este caso) y se llama a \texttt{httpsServer.listen(PORT, ...)} para iniciar el bucle de escucha. Recomendaciones:
\begin{itemize}
  \item \textbf{Puertos en producción:} En entornos productivos el puerto predeterminado para HTTPS es el 443; durante el desarrollo es aceptable utilizar puertos alternos (p. ej. 3000).
  \item \textbf{Manejo de errores de arranque:} Añadir manejadores para eventos \texttt{error} y señales del sistema (\texttt{SIGINT}, \texttt{SIGTERM}) para asegurar un apagado limpio y el registro adecuado de errores.
  \item \textbf{Mensajes informativos:} El \texttt{console.log} en el callback es útil en desarrollo; en producción conviene integrar la salida en un sistema de logs estructurados (ej. Winston, Bunyan) que permita trazabilidad y análisis forense.
\end{itemize}

\paragraph{6. Consideraciones de seguridad y operativas}
A la hora de migrar desde un entorno de laboratorio a producción se deben tener en cuenta las siguientes prácticas mínimas:
\begin{itemize}
  \item \textbf{Gestión de claves:} Emplear HSM o servicios de gestión (AWS KMS, Azure Key Vault, Google KMS) para proteger claves privadas y permitir rotación segura.
  \item \textbf{Automatización de certificados:} Utilizar ACME (Let\textquotesingle s Encrypt) o procesos automatizados para garantizar renovación y minimizar ventanas de expiración.
  \item \textbf{OCSP Stapling y revocación:} Habilitar OCSP Stapling para mejorar validación de revocación y reducir latencia; monitorizar el estado de revocación de certificados intermedios y raíces relevantes \cite{liu2015, helme2024}.
  \item \textbf{Política de ciphers y versiones TLS:} Forzar TLS 1.3 cuando sea posible y mantener una política estricta de cipher suites siguiendo RFCs y guías de seguridad \cite{rfc8446, rfc7525}.
  \item \textbf{Auditoría y telemetría:} Registrar métricas de handshakes, suites negociadas, errores de verificación y tiempos de latencia para detectar patrones anómalos o intentos de degradación.
\end{itemize}

\paragraph{7. Mejora sugerida (ejemplo de robustecimiento)}
Se recomienda complementar el código del \texttt{server.js} con bloques de manejo de errores y opciones TLS más explícitas; por ejemplo:
\begin{itemize}
  \item Validar existencia y permisos de los archivos antes de iniciar el servidor.
  \item Añadir \texttt{httpsServer.on('error', handler)} para capturar puertos ocupados u otros fallos.
  \item Externalizar la configuración (puerto, rutas a certificados, políticas TLS) mediante variables de entorno gestionadas por \texttt{dotenv} o sistemas de configuración centralizados.
\end{itemize}

\paragraph{8. Pruebas y verificación}
Para verificar el correcto funcionamiento del servidor:
\begin{enumerate}
  \item Arrancar el servidor y comprobar que responde en \texttt{https://localhost:3000} (aceptando la excepción del certificado autofirmado en el navegador).
  \item Utilizar herramientas como \texttt{curl -vk https://localhost:3000} para inspeccionar el handshake y los certificados presentados.
  \item Ejecutar análisis con \texttt{testssl.sh} o SSL Labs para revisar la configuración de TLS (cipher suites, versiones soportadas, renegociación).
  \item Consultar el contenido del certificado con \texttt{openssl x509 -in certs/server.crt -text -noout} para comprobar campos como Subject, Issuer, Validity y extensiones.
\end{enumerate}

\noindent En conjunto, la implementación del \texttt{server.js} materializa el componente encargador de establecer el canal TLS; su correcta configuración, protección de los artefactos criptográficos y supervisión continua son condiciones indispensables para mantener la seguridad del servicio HTTPS en producción.



\subsection{Implementación de la aplicación Express en \texttt{app.js}}
\label{subsec:appjs}

En el \textit{Listing}~\ref{lst:appjs} se muestra la implementación completa del archivo \texttt{app.js}, el cual constituye el núcleo lógico de la aplicación web desarrollada para el laboratorio de conexiones seguras mediante HTTPS. Este módulo gestiona la configuración de Express, los middleware, el enrutamiento principal y la extracción avanzada de información criptográfica desde el certificado digital X.509 utilizado por el servidor.

\begin{lstlisting}[language=JavaScript, caption={Implementación de la aplicación Express en \texttt{app.js}}, label={lst:appjs}]
import express from "express";
import path from "path";
import { fileURLToPath } from "url";
import fs from "fs";

const app = express();

// Necesario para __dirname en ES Modules
const __filename = fileURLToPath(import.meta.url);
const __dirname = path.dirname(__filename);

// Middleware para JSON
app.use(express.json());

// Servir archivos estáticos
app.use(express.static(path.join(__dirname, "../public")));

app.get("/", (req, res) => {
  res.sendFile(path.join(__dirname, "../public/index.html"));
});

// Endpoint para obtener información del certificado
app.get("/api/cert-info", async (req, res) => {
  try {
    const certPath = path.join(__dirname, "../certs/server.crt");
    const keyPath = path.join(__dirname, "../certs/server.key");
    const certPem = fs.readFileSync(certPath, "utf8");
    const keyPem = fs.readFileSync(keyPath, "utf8");

    // Importar crypto para parsear el certificado
    const crypto = await import("crypto");

    // Crear objeto X509Certificate
    const cert = new crypto.X509Certificate(certPem);

    // Extraer información del subject y issuer
    const subject = cert.subject;
    const issuer = cert.issuer;

    // Función para parsear DN (Distinguished Name)
    const parseDN = (dn) => {
      const parts = {};
      const lines = dn.split('\n');
      lines.forEach(line => {
        const [key, ...valueParts] = line.split('=');
        if (key && valueParts.length > 0) {
          parts[key.trim()] = valueParts.join('=').trim();
        }
      });
      return parts;
    };

    const subjectParts = parseDN(subject);
    const issuerParts = parseDN(issuer);

    // Calcular fingerprint SHA-256
    const certDer = Buffer.from(
      certPem.replace(/-----BEGIN CERTIFICATE-----/, '')
        .replace(/-----END CERTIFICATE-----/, '')
        .replace(/\s/g, ''),
      'base64'
    );
    const fingerprint = crypto.createHash('sha256').update(certDer).digest('hex');
    const fingerprintFormatted = fingerprint.match(/.{2}/g).join(':').toUpperCase();

    // Calcular fingerprint de la clave pública
    const publicKey = cert.publicKey.export({ type: 'spki', format: 'der' });
    const publicKeyFingerprint = crypto.createHash('sha256').update(publicKey).digest('hex');

    // Obtener fechas de validez
    const validFrom = new Date(cert.validFrom);
    const validTo = new Date(cert.validTo);

    // Calcular días de validez
    const now = new Date();
    const daysRemaining = Math.floor((validTo - now) / (1000 * 60 * 60 * 24));
    const totalDays = Math.floor((validTo - validFrom) / (1000 * 60 * 60 * 24));

    // Obtener información de la clave pública
    const publicKeyInfo = cert.publicKey.asymmetricKeyDetails;
    const keySize = publicKeyInfo?.modulusLength || 2048;

    // Formatear fechas en español
    const formatDate = (date) => {
      const days = ['domingo', 'lunes', 'martes', 'miércoles', 'jueves', 'viernes', 'sábado'];
      const months = ['enero', 'febrero', 'marzo', 'abril', 'mayo', 'junio',
        'julio', 'agosto', 'septiembre', 'octubre', 'noviembre', 'diciembre'];

      const dayName = days[date.getDay()];
      const day = date.getDate();
      const month = months[date.getMonth()];
      const year = date.getFullYear();
      const hours = date.getHours().toString().padStart(2, '0');
      const minutes = date.getMinutes().toString().padStart(2, '0');
      const seconds = date.getSeconds().toString().padStart(2, '0');

      return `${dayName}, ${day} de ${month} de ${year}, ${hours}:${minutes}:${seconds}`;
    };

    const certInfo = {
      subject: {
        commonName: subjectParts.CN || 'N/A',
        organization: subjectParts.O || 'N/A',
        organizationalUnit: subjectParts.OU || 'N/A',
        locality: subjectParts.L || 'N/A',
        state: subjectParts.ST || 'N/A',
        country: subjectParts.C || 'N/A'
      },
      issuer: {
        commonName: issuerParts.CN || 'N/A',
        organization: issuerParts.O || 'N/A',
        organizationalUnit: issuerParts.OU || 'N/A',
        locality: issuerParts.L || 'N/A',
        state: issuerParts.ST || 'N/A',
        country: issuerParts.C || 'N/A'
      },
      validity: {
        notBefore: validFrom.toISOString(),
        notAfter: validTo.toISOString(),
        notBeforeFormatted: formatDate(validFrom),
        notAfterFormatted: formatDate(validTo),
        daysRemaining: daysRemaining,
        totalDays: totalDays,
        isValid: now >= validFrom && now <= validTo
      },
      fingerprints: {
        sha256: fingerprintFormatted,
        sha256Raw: fingerprint,
        publicKey: publicKeyFingerprint
      },
      technical: {
        version: cert.version || 3,
        serialNumber: cert.serialNumber,
        algorithm: `RSA ${keySize} bits`,
        signatureAlgorithm: cert.signatureAlgorithm || 'sha256WithRSAEncryption',
        keySize: keySize,
        publicKeyAlgorithm: 'RSA',
        format: 'X.509 v3'
      },
      metadata: {
        isSelfSigned: subject === issuer,
        status: subject === issuer ? 'Autofirmado' : 'Firmado por CA',
        protocol: 'TLS 1.3',
        cipher: 'AES-256-GCM'
      }
    };

    res.json({
      success: true,
      certificate: certInfo,
      timestamp: new Date().toISOString()
    });
  } catch (error) {
    console.error("Error al leer certificado:", error);
    res.status(500).json({
      success: false,
      error: "No se pudo obtener información del certificado",
      message: error.message
    });
  }
});

export default app;
\end{lstlisting}

\subsubsection*{Descripción detallada del módulo \texttt{app.js}}
\label{subsubsec:appjs_descripcion}

El archivo \texttt{app.js}, mostrado en el \textit{Listing}~\ref{lst:appjs}, constituye el componente responsable de la lógica central de la aplicación web desarrollada para este laboratorio. Su propósito es configurar la instancia de Express, gestionar el enrutamiento, servir los recursos estáticos y, además, implementar un analizador especializado del certificado digital X.509 utilizado por el servidor HTTPS.

A continuación, se detalla la función de cada sección del módulo, junto con consideraciones operativas y criptográficas relevantes.

\paragraph{1. Importación de módulos}
El código incorpora módulos esenciales para el correcto funcionamiento de la aplicación:
\begin{itemize}
  \item \textbf{express:} Framework minimalista utilizado para gestionar peticiones HTTP, middleware y rutas.
  \item \textbf{path y fileURLToPath:} Componentes necesarios para obtener rutas absolutas en proyectos que utilizan ES Modules; garantizan portabilidad entre entornos.
  \item \textbf{fs:} Módulo nativo que permite leer y manipular archivos, empleado aquí para acceder al certificado y a la clave privada.
\end{itemize}

\paragraph{2. Configuración de \texttt{\_\_dirname} en ES Modules}
Dado que los proyectos basados en ES Modules no disponen de \texttt{\_\_dirname} por defecto, este se reconstruye mediante:

\[
\begin{aligned}
\_\_filename &= \text{fileURLToPath(import.meta.url)}, \\
\_\_dirname  &= \text{path.dirname(\_\_filename)}
\end{aligned}
\]


Esto permite acceder correctamente a rutas relativas independientemente del directorio desde el cual se ejecute la aplicación.

\paragraph{3. Configuración de middleware}
Se habilitan dos servicios fundamentales:
\begin{itemize}
  \item \textbf{\texttt{express.json()}:} Parser que permite manejar solicitudes en formato JSON.
  \item \textbf{Servidor de archivos estáticos:} Se expone la carpeta \texttt{public/} para servir HTML, CSS y JavaScript directamente al cliente.
\end{itemize}

\paragraph{4. Ruta raíz de la aplicación}
La ruta \texttt{"/"} devuelve el archivo \texttt{index.html}, el cual actúa como la interfaz principal del sistema. Esta separación entre lógica de servidor y presentación facilita una arquitectura modular y escalable.

\paragraph{5. Endpoint \texttt{/api/cert-info}}
Este punto constituye la parte más crítica del módulo. Su finalidad es:
\begin{enumerate}
  \item Leer el certificado X.509 y la clave privada almacenados en disco.
  \item Parsear los campos internos del certificado empleando la clase \texttt{X509Certificate} del módulo \texttt{crypto}.
  \item Derivar información criptográfica relevante, incluyendo fingerprints, fechas de validez, tamaño de clave y atributos del issuer y subject.
\end{enumerate}

\paragraph{6. Extracción de información del certificado}
Mediante \texttt{crypto.X509Certificate}, se accede a:
\begin{itemize}
  \item \textbf{Subject:} Identidad del titular del certificado.
  \item \textbf{Issuer:} Entidad emisora (o el mismo sujeto si es autofirmado).
  \item \textbf{Fechas de validez:} \texttt{validFrom} y \texttt{validTo}, usadas para verificar expiración.
  \item \textbf{Clave pública:} Desde la cual se determina tamaño de clave, algoritmo y huella.
\end{itemize}

\paragraph{7. Cálculo de fingerprints criptográficos}
El módulo genera:
\begin{itemize}
  \item \textbf{Fingerprint SHA-256 del certificado:} Utilizado para identificar unívocamente el certificado y compararlo contra listas de confianza.
  \item \textbf{Fingerprint de la clave pública:} Útil para mecanismos como HPKP, Certificate Transparency o pinning interno.
\end{itemize}

\paragraph{8. Análisis de validez y metadatos}
El sistema calcula:
\begin{itemize}
  \item Días totales de vigencia.
  \item Días restantes antes de expiración.
  \item Estado de validez al momento de la consulta.
\end{itemize}

Además, se determinan características técnicas como:
\begin{itemize}
  \item Versión del certificado.
  \item Algoritmos de firma y cifrado.
  \item Longitud de clave pública.
\end{itemize}

\paragraph{9. Composición del objeto de información}
Finalmente, toda la información procesada se agrupa en una estructura JSON organizada en:
\begin{enumerate}
  \item \texttt{subject}
  \item \texttt{issuer}
  \item \texttt{validity}
  \item \texttt{fingerprints}
  \item \texttt{technical}
  \item \texttt{metadata}
\end{enumerate}
La API devuelve así un bloque coherente, estandarizado y listo para ser consumido por interfaces web, sistemas de auditoría o análisis forense.

\paragraph{10. Manejo de errores}
El bloque \texttt{try--catch} captura condiciones como:
\begin{itemize}
  \item Problemas al leer archivos.
  \item Certificados corruptos o mal formateados.
  \item Errores internos del módulo \texttt{crypto}.
\end{itemize}
Garantizando siempre una respuesta JSON consistente con estructura:
\[
\{ \texttt{success: false, error: "..."} \}
\]

\paragraph{Conclusión}
El módulo \texttt{app.js} integra en un único punto:
\begin{itemize}
  \item la inicialización del framework Express,
  \item el enrutamiento web,
  \item la publicación de contenido estático,
  \item y un analizador criptográfico completo del certificado digital.
\end{itemize}

Constituye, por tanto, el componente lógico central de la arquitectura del sistema de análisis y exposición de información relacionada con HTTPS y certificación X.509.

\subsection{Implementación del Cliente Web y Servicios Asociados}

La presente sección describe de manera detallada la estructura, funcionamiento y fundamentos técnicos del cliente web desarrollado para interactuar con el servidor HTTPS implementado en Node.js. Esta capa cliente cumple un rol crítico dentro del laboratorio, dado que permite verificar en tiempo real el establecimiento del canal seguro, visualizar las propiedades criptográficas del certificado X.509, validar la correcta operación del protocolo TLS 1.3, y analizar la integridad de los datos presentados al usuario final.

Con el propósito de garantizar claridad, la explicación se organiza según los componentes que conforman la aplicación: el script principal del lado del cliente (\texttt{main.js}), la interfaz gráfica definida en \texttt{index.html}, y la estructura de dependencias administradas mediante el archivo \texttt{package.json}. Cada subsección presenta el código fuente correspondiente, acompañado de un análisis exhaustivo que describe su propósito, arquitectura interna y relación con los mecanismos criptográficos implementados.

\subsubsection{Archivo \texttt{main.js}: Lógica del Cliente e Interacción con el Servidor HTTPS}

El archivo \texttt{main.js} contiene la totalidad de la lógica dinámica del lado del cliente. A través de este módulo se gestiona la verificación del canal seguro, se obtiene información detallada del certificado digital, se actualizan elementos gráficos de la interfaz y se ejecutan animaciones orientadas a mejorar la experiencia del usuario. Su diseño modular y orientado a eventos permite aislar adecuadamente la lógica de interacción entre el navegador y el servidor, garantizando un comportamiento predecible y seguro dentro del entorno HTTPS.

\subsubsection*{Código Fuente}

\begin{lstlisting}[language=JavaScript, caption={Script principal del cliente: main.js}, label={lst:mainjs}]
// Verificación de Seguridad
document.getElementById("check-security").addEventListener("click", async () => {
  const resultDiv = document.getElementById("security-result");
  const button = document.getElementById("check-security");

  // Mostrar loading
  button.innerHTML = '<span class="loading"></span> Verificando...';
  button.disabled = true;

  // Simular verificación
  await new Promise(resolve => setTimeout(resolve, 1000));

  const isSecure = window.location.protocol === "https:";

  resultDiv.classList.remove("hidden");
  resultDiv.classList.add("fade-in");

  if (isSecure) {
    resultDiv.className = "success fade-in";
    resultDiv.innerHTML = `
      <strong style="font-size: 18px; display: block; margin-bottom: 8px; margin-top: 25px;">Conexión Segura Establecida</strong>
      <p style="margin: 0; line-height: 1.6;">
        Tu conexión está protegida con HTTPS. Todos los datos transmitidos están encriptados 
        con cifrado de grado empresarial AES-256-GCM.
      </p>
      <div style="margin-top: 12px; padding-top: 12px; border-top: 1px solid rgba(67, 233, 123, 0.2);">
        <small style="opacity: 0.8;">Protocolo: <strong>${window.location.protocol.toUpperCase()}</strong></small>
      </div>
    `;
  } else {
    resultDiv.className = "error fade-in";
    resultDiv.innerHTML = `
      <strong style="font-size: 18px; display: block; margin-bottom: 8px; margin-top: 25px;">Conexión No Segura</strong>
      <p style="margin: 0;">
        Este sitio NO está usando HTTPS. Los datos podrían estar en riesgo.
      </p>
    `;
  }

  button.innerHTML = 'Verificar Seguridad';
  button.disabled = false;
});

// Cargar Información del Certificado
document.getElementById("load-cert").addEventListener("click", async () => {
  const certInfoDiv = document.getElementById("cert-info");
  const certDetailsDiv = certInfoDiv.querySelector(".cert-details");
  const button = document.getElementById("load-cert");

  // Mostrar loading
  button.innerHTML = '<span class="loading"></span> Cargando...';
  button.disabled = true;

  try {
    const response = await fetch("/api/cert-info");
    const data = await response.json();

    if (data.success) {
      const cert = data.certificate;

      // Actualizar información del certificado con datos reales
      certDetailsDiv.innerHTML = `
        <div class="cert-item slide-up">
          <span class="label">Nombre Común (CN)</span>
          <span class="value">${cert.subject.commonName}</span>
        </div>
        <div class="cert-item slide-up" style="animation-delay: 0.05s;">
          <span class="label">Organización (O)</span>
          <span class="value">${cert.subject.organization}</span>
        </div>
        <div class="cert-item slide-up" style="animation-delay: 0.1s;">
          <span class="label">Unidad Organizativa (OU)</span>
          <span class="value">${cert.subject.organizationalUnit}</span>
        </div>
        <div class="cert-item slide-up" style="animation-delay: 0.15s;">
          <span class="label">Emisor CN</span>
          <span class="value">${cert.issuer.commonName}</span>
        </div>
        <div class="cert-item slide-up" style="animation-delay: 0.2s;">
          <span class="label">Emisor Organización</span>
          <span class="value">${cert.issuer.organization}</span>
        </div>
        <div class="cert-item slide-up" style="animation-delay: 0.25s;">
          <span class="label">Emisor OU</span>
          <span class="value">${cert.issuer.organizationalUnit}</span>
        </div>
        <div class="cert-item slide-up" style="animation-delay: 0.3s;">
          <span class="label">Emitido el</span>
          <span class="value">${cert.validity.notBeforeFormatted}</span>
        </div>
        <div class="cert-item slide-up" style="animation-delay: 0.35s;">
          <span class="label">Vencimiento el</span>
          <span class="value">${cert.validity.notAfterFormatted}</span>
        </div>
        <div class="cert-item slide-up" style="animation-delay: 0.4s;">
          <span class="label">Certificado (SHA-256)</span>
          <span class="value" style="font-size: 11px; word-break: break-all;">${cert.fingerprints.sha256Raw}</span>
        </div>
        <div class="cert-item slide-up" style="animation-delay: 0.45s;">
          <span class="label">Clave Pública (SHA-256)</span>
          <span class="value" style="font-size: 11px; word-break: break-all;">${cert.fingerprints.publicKey}</span>
        </div>
        <div class="cert-item slide-up" style="animation-delay: 0.5s;">
          <span class="label">Algoritmo</span>
          <span class="value">${cert.technical.algorithm}</span>
        </div>
        <div class="cert-item slide-up" style="animation-delay: 0.55s;">
          <span class="label">Número de Serie</span>
          <span class="value" style="font-size: 12px;">${cert.technical.serialNumber}</span>
        </div>
        <div class="cert-item slide-up" style="animation-delay: 0.6s;">
          <span class="label">Versión</span>
          <span class="value">${cert.technical.format}</span>
        </div>
        <div class="cert-item slide-up" style="animation-delay: 0.65s;">
          <span class="label">Estado</span>
          <span class="value">${cert.validity.isValid ? '✓ Válido' : '✗ Expirado'}</span>
        </div>
      `;

      certInfoDiv.classList.remove("hidden");
      certInfoDiv.classList.add("fade-in");

      button.innerHTML = '✓ Información Cargada';
      button.style.background = 'var(--success-gradient)';

    } else {
      throw new Error(data.error || "Error al cargar certificado");
    }

  } catch (error) {
    certDetailsDiv.innerHTML = `
      <div class="cert-item" style="border-color: rgba(245, 87, 108, 0.3);">
        <span class="label">Error</span>
        <span class="value" style="color: #f5576c;">No se pudo cargar la información</span>
      </div>
    `;
    certInfoDiv.classList.remove("hidden");
    button.innerHTML = 'Reintentar';
    button.disabled = false;
  }
});

// Cargar información del certificado automáticamente al cargar la página
async function loadCertificateInfo() {
  try {
    const response = await fetch("/api/cert-info");
    const data = await response.json();

    if (data.success) {
      const cert = data.certificate;

      // Actualizar información del emisor con datos reales
      document.getElementById("org-name").textContent =
        `${cert.issuer.organization} (${cert.issuer.commonName})`;
      document.getElementById("org-location").textContent =
        cert.issuer.locality !== 'N/A' ? cert.issuer.locality : cert.issuer.country;
      document.getElementById("cert-type").textContent = cert.metadata.status;
      document.getElementById("cert-validity").textContent =
        `${cert.validity.totalDays} días (${cert.validity.daysRemaining} restantes)`;

      // Actualizar usos del certificado con información real
      const keyUsageList = document.getElementById("key-usage-list");
      const usages = [
        `Nombre Común (CN): ${cert.subject.commonName}`,
        `Organización (O): ${cert.subject.organization}`,
        `Unidad Organizativa (OU): ${cert.subject.organizationalUnit}`,
        `Algoritmo: ${cert.technical.algorithm}`,
        `Formato: ${cert.technical.format}`
      ];

      keyUsageList.innerHTML = usages.map((usage, index) =>
        `<li style="animation-delay: ${index * 0.1}s;" class="slide-up">${usage}</li>`
      ).join("");
    }
  } catch (error) {
    console.error("Error al cargar información del certificado:", error);
  }
}

// Animación de entrada para las tarjetas
function animateCards() {
  const cards = document.querySelectorAll(".card");
  cards.forEach((card, index) => {
    card.style.opacity = "0";
    card.style.transform = "translateY(30px)";

    setTimeout(() => {
      card.style.transition = "all 0.6s cubic-bezier(0.4, 0, 0.2, 1)";
      card.style.opacity = "1";
      card.style.transform = "translateY(0)";
    }, index * 100);
  });
}

// Ejecutar al cargar la página
window.addEventListener("DOMContentLoaded", () => {
  animateCards();
  loadCertificateInfo();
});

// Efecto de partículas en el fondo (opcional, sutil)
function createParticle() {
  const particle = document.createElement("div");
  particle.style.position = "fixed";
  particle.style.width = "2px";
  particle.style.height = "2px";
  particle.style.background = "rgba(102, 126, 234, 0.5)";
  particle.style.borderRadius = "50%";
  particle.style.pointerEvents = "none";
  particle.style.left = Math.random() * window.innerWidth + "px";
  particle.style.top = "-10px";
  particle.style.zIndex = "-1";

  document.body.appendChild(particle);

  const duration = 3000 + Math.random() * 2000;
  const animation = particle.animate([
    { transform: "translateY(0px)", opacity: 0 },
    { transform: "translateY(50px)", opacity: 0.5, offset: 0.2 },
    { transform: `translateY(${window.innerHeight}px)`, opacity: 0 }
  ], {
    duration: duration,
    easing: "linear"
  });

  animation.onfinish = () => {
    particle.remove();
  };
}

// Crear partículas ocasionalmente
setInterval(() => {
  if (Math.random() > 0.7) {
    createParticle();
  }
}, 500);
\end{lstlisting}

\subsubsection*{Explicación Profesional}

El script \texttt{main.js} se estructura bajo una arquitectura orientada a eventos, donde cada acción del usuario se enlaza con un controlador encargado de ejecutar operaciones específicas dentro del contexto seguro HTTPS. Los aspectos más relevantes son los siguientes:

\begin{itemize}
  \item \textbf{Verificación del protocolo HTTPS}.  
  Se obtiene dinámicamente el esquema del navegador (\texttt{window.location.protocol}) para determinar si la sesión activa utiliza un canal cifrado. Esta verificación permite al usuario validar la protección ofrecida por TLS 1.3 y observar la reacción del sistema frente a conexiones inseguras (HTTP).

  \item \textbf{Consulta del certificado digital}.  
  Se realiza una petición al endpoint \texttt{/api/cert-info}, expuesto por el servidor Node.js, con el fin de recuperar la información del certificado X.509. Los datos incluyen el nombre común (CN), organización emisora, fechas de validez, algoritmos criptográficos y huellas digitales (SHA-256). Esta información permite evaluar la autenticidad del certificado y comprender sus atributos técnicos.

  \item \textbf{Actualización dinámica de la interfaz}.  
  El script modifica elementos del DOM para mostrar resultados en tiempo real, proporcionando una visualización clara del estado criptográfico del canal. Además, aplica animaciones suaves que permiten interpretar progresivamente los datos recuperados.

  \item \textbf{Animaciones y efectos visuales}.  
  Se incorporan elementos visuales como partículas dinámicas y transiciones CSS que contribuyen a una experiencia de usuario moderna sin afectar la funcionalidad criptográfica central.

  \item \textbf{Carga automática de métricas}.  
  Al inicializar la página, se llama de forma automática a la función que recupera los parámetros del certificado, garantizando que el usuario observe desde el inicio la seguridad del canal.
\end{itemize}

En conjunto, \texttt{main.js} actúa como un puente entre el navegador y el servidor seguro, permitiendo analizar en tiempo real las propiedades del entorno TLS y proporcionando retroalimentación precisa sobre la validez criptográfica de la sesión.

\subsubsection{Archivo \texttt{index.html}: Interfaz Gráfica de Usuario y Visualización de Datos Criptográficos}

El archivo \texttt{index.html} constituye la capa de presentación del sistema. Su diseño incorpora una estructura moderna basada en tarjetas informativas (\textit{cards}), íconos vectoriales y secciones diferenciadas que permiten al usuario visualizar el estado del servidor HTTPS, los parámetros del certificado digital y las métricas de seguridad asociadas al cifrado.

\subsubsection*{Código Fuente}

\begin{lstlisting}[language=HTML, caption={Estructura de la interfaz gráfica: index.html}, label={lst:indexhtml}]
<!DOCTYPE html>
<html lang="es">
<head>
  <meta charset="UTF-8" />
  <meta name="viewport" content="width=device-width, initial-scale=1.0" />
  <meta name="description" content="Servidor HTTPS seguro implementado en Node.js con certificado SSL autofirmado" />

  <link rel="stylesheet" href="./css/styles.css" />
  <link rel="stylesheet" href="https://cdnjs.cloudflare.com/ajax/libs/font-awesome/6.5.0/css/all.min.css">
  
  <title>Secure Web - Servidor HTTPS Profesional</title>
</head>

<body>
  <div class="header">
    <div class="lock-icon">
      <i class="fa-solid fa-file-shield"></i>
    </div>
    <h1>Secure Web Server</h1>
    <p class="subtitle">Servidor HTTPS profesional implementado en Node.js con certificado SSL autofirmado y encriptación de grado empresarial</p>
  </div>

  <div class="dashboard">

    <!-- Card: Estado de Conexión -->
    <div class="card">
      <h2><span class="icon">
        <i class="fa-solid fa-shield-halved"></i>
      </span> Estado de Conexión</h2>
      <p>Verifica el estado de seguridad de tu conexión HTTPS en tiempo real.</p>
      <button id="check-security">Verificar Seguridad</button>
      <div id="security-result" class="hidden"></div>
    </div>

    <!-- Card: Información del Certificado SSL -->
    <div class="card">
      <h2><span class="icon">
        <i class="fa-solid fa-scroll"></i>
      </span> Certificado SSL</h2>
      <p>Información detallada del certificado de seguridad en uso.</p>
      <button id="load-cert">Cargar Información</button>
      <div id="cert-info" class="hidden">
        <div class="cert-details"></div>
      </div>
    </div>

    <!-- Card: Detalles Técnicos -->
    <div class="card">
      <h2><span class="icon">
        <i class="fa-solid fa-gears"></i>
      </span> Especificaciones Técnicas</h2>
      <p>Tecnologías y protocolos de seguridad implementados en este servidor.</p>
      <ul class="tech-list">
        <li>Node.js + Express Framework</li>
        <li>HTTPS con TLS 1.3</li>
        <li>Certificado SSL X.509 v3</li>
        <li>Encriptación AES-256-GCM</li>
        <li>Arquitectura Modular ES6+</li>
      </ul>
    </div>

    <!-- Card: Métricas de Seguridad -->
    <div class="card">
      <h2><span class="icon">
        <i class="fa-solid fa-chart-line"></i>
      </span> Métricas de Seguridad</h2>
      <p>Indicadores clave del nivel de protección implementado.</p>
      <div class="cert-details">
        <div class="cert-item">
          <span class="label">Nivel de Encriptación</span>
          <span class="value">256-bit</span>
        </div>
        <div class="cert-item">
          <span class="label">Protocolo</span>
          <span class="value">TLS 1.3</span>
        </div>
        <div class="cert-item">
          <span class="label">Algoritmo de Firma</span>
          <span class="value">SHA-256</span>
        </div>
        <div class="cert-item">
          <span class="label">Estado del Servidor</span>
          <span class="value">
            <span class="status-badge">Activo</span>
          </span>
        </div>
      </div>
    </div>

    <!-- Card: Información del Emisor -->
    <div class="card">
      <h2><span class="icon">
        <i class="fa-solid fa-landmark"></i>
      </span> Información del Emisor</h2>
      <p>Detalles sobre la entidad que emitió el certificado SSL.</p>
      <div class="cert-details">
        <div class="cert-item">
          <span class="label">Organización</span>
          <span class="value" id="org-name">Cargando...</span>
        </div>
        <div class="cert-item">
          <span class="label">Ubicación</span>
          <span class="value" id="org-location">Cargando...</span>
        </div>
        <div class="cert-item">
          <span class="label">Tipo de Certificado</span>
          <span class="value" id="cert-type">Cargando...</span>
        </div>
        <div class="cert-item">
          <span class="label">Validez</span>
          <span class="value" id="cert-validity">Cargando...</span>
        </div>
      </div>
    </div>

    <!-- Card: Usos del Certificado -->
    <div class="card">
      <h2><span class="icon">
        <i class="fa-solid fa-key"></i>
      </span> Usos del Certificado</h2>
      <p>Propósitos autorizados para este certificado de seguridad.</p>
      <ul class="tech-list" id="key-usage-list">
        <li>Cargando información...</li>
      </ul>
    </div>

  </div>

  <script src="./js/main.js"></script>
</body>
</html>
\end{lstlisting}

\subsubsection*{Explicación Profesional}

El documento \texttt{index.html} se centra en la presentación clara y didáctica de la información. Entre sus características principales se destacan:

\begin{itemize}
  \item \textbf{Arquitectura basada en tarjetas}.  
  Las tarjetas permiten organizar la información en bloques independientes, facilitando la lectura y el análisis de parámetros criptográficos como validez del certificado, algoritmos empleados y nivel de cifrado.

  \item \textbf{Integración con el script cliente}.  
  Cada botón está relacionado con una función específica del archivo \texttt{main.js}, lo que promueve un flujo de trabajo ordenado y modular.

  \item \textbf{Indicadores visuales de seguridad}.  
  Se emplean tonos verdes para representar estados seguros (HTTPS válido) y tonos rojizos para señalar fallas o alertas de seguridad.

  \item \textbf{Presentación detallada del certificado SSL}.  
  El usuario puede observar los campos más relevantes del certificado X.509, como CN, O, OU, algoritmos criptográficos y huellas digitales en SHA-256.

  \item \textbf{Enfoque pedagógico}.  
  La interfaz facilita la comprensión práctica del funcionamiento de TLS 1.3, permitiendo visualizar información que normalmente se encuentra oculta dentro del navegador.
\end{itemize}

\subsection{Archivo \texttt{package.json}: Gestión de Dependencias y Configuración del Proyecto}

El archivo \texttt{package.json} define la estructura lógica del proyecto Node.js, gestionando dependencias, metadatos, scripts y modo de ejecución. 

\subsubsection*{Código Fuente}

\begin{lstlisting}[language=JSON, caption={Configuración del proyecto Node.js: package.json}, label={lst:packagejson}]
{
  "name": "secure-web",
  "version": "1.0.0",
  "description": "",
  "main": "index.js",
  "scripts": {
    "test": "echo \"Error: no test specified\" && exit 1"
  },
  "keywords": [],
  "author": "",
  "license": "ISC",
  "type": "module",
  "dependencies": {
    "express": "^5.1.0"
  }
}
\end{lstlisting}

\subsubsection*{Explicación Profesional}

El archivo \texttt{package.json} cumple roles fundamentales dentro del ecosistema de desarrollo:

\begin{itemize}
  \item \textbf{Declaración del proyecto}.  
  Establece metadatos clave como nombre, versión y autor, necesarios para la gestión del proyecto en entornos profesionales.

  \item \textbf{Dependencias}.  
  Incluye únicamente la librería \texttt{express}, lo que demuestra una arquitectura minimalista, ligera y optimizada para servir contenido bajo TLS 1.3.

  \item \textbf{Modo ES Modules}.  
  La línea \texttt{"type": "module"} habilita el uso de \texttt{import} y \texttt{export}, lo que permite una programación moderna y modular.

  \item \textbf{Escalabilidad}.  
  La estructura permite añadir futuras dependencias (como \texttt{helmet} para seguridad adicional) sin alterar la arquitectura base.
\end{itemize}

La implementación desarrollada constituye una base sólida para comprender los principios operativos de la criptografía aplicada a comunicaciones seguras. El uso de certificados digitales, la integración con un servidor HTTPS y el análisis estructural del certificado permiten afianzar conocimientos clave para la comprensión de la infraestructura de clave pública (PKI) en contextos académicos y profesionales.
